\clearpage
\subsubsection{Section Reproduction Steps}\label{subsubsec:reproduce-financial}
From the repository root, after \texttt{make setup}, replay the retained May--July
2026 extracts and rebuild the PDF:
\begin{lstlisting}[basicstyle=\ttfamily\small,breaklines=true]
make reproduce-financial-section
make report
\end{lstlisting}
The section target uses cached inputs only. It recreates
Table~\ref{tab:financial-opportunity-cancellations}, split by cancellation reason
and line-item status, with
\path{opportunity-order-cancellation-summary} using
\path{DATE_BASIS=shipped-or-cancelled} and \texttt{WRITE\_LATEX=1}.
The four follow-up comparisons use
\path{opportunity-follow-up-sales},
\path{opportunity-new-upload-customer-follow-up} and
\path{opportunity-new-vet-diet-customer-follow-up}, including their Chewy-wide
companions, then \path{combined-follow-up-table}.
\path{financial-sql-listings} recreates the executable SQL listings.
All use the displayed May 1 inclusive--August 1 exclusive cohort and October 28,
2026 observation cutoff; follow-up windows are 7 through 70 days.

Retained Parquet extracts, manifests and first-order cohort inputs under
\path{data/opportunity/} are required and are not stored in Git.
Fingerprints are checked before reuse. To obtain a new warehouse snapshot,
run the underlying targets without \texttt{CACHED=1}, following HOWTO.md for the
first-order history prerequisites. New snapshots may change counts.
